% paper12/matrix_spencer_lifts_DRAFT.tex
% RECEIVED DRAFT (owner-supplied, GPT-assisted rewrite) — to be VERIFIED and REPAIRED.
% Known mechanical issues to fix in the clean version (matrix_spencer_lifts.tex):
%   (1) Display-math delimiters arrived as bare [ ... ] (backslashes stripped in transit) -> restore \[ ... \].
%   (2) Every theorem/lemma/prop uses an identical \label{thm}; all \ref{thm} are therefore broken -> uniquify.
%   (3) \DeclareMathOperator*{\argmin}{arg,min} should be arg\,min.
% Mathematical content is preserved verbatim for the verification lane to check.
%
\documentclass[11pt]{amsart}

\usepackage[T1]{fontenc}
\usepackage[utf8]{inputenc}
\usepackage{lmodern}
\usepackage{microtype}
\usepackage{amsmath,amssymb,amsthm,mathtools}
\usepackage{mathrsfs}
\usepackage{enumitem}
\usepackage[colorlinks=true,linkcolor=blue,citecolor=blue,urlcolor=blue]{hyperref}

\numberwithin{equation}{section}

\newtheorem{theorem}{Theorem}[section]
\newtheorem{proposition}[theorem]{Proposition}
\newtheorem{corollary}[theorem]{Corollary}
\newtheorem{lemma}[theorem]{Lemma}
\newtheorem{question}[theorem]{Question}
\theoremstyle{definition}
\newtheorem{definition}[theorem]{Definition}
\theoremstyle{remark}
\newtheorem{remark}[theorem]{Remark}

\newcommand{\C}{\mathbb C}
\newcommand{\R}{\mathbb R}
\newcommand{\Z}{\mathbb Z}
\newcommand{\T}{\mathbb T}
\newcommand{\M}{\mathrm M}
\newcommand{\Id}{I}
\newcommand{\op}{\mathrm{op}}
\newcommand{\Span}{\operatorname{span}}
\newcommand{\disc}{\operatorname{disc}}
\newcommand{\discop}{\operatorname{disc}_{\op}}
\newcommand{\Ent}{\mathsf E}
\newcommand{\supp}{\operatorname{supp}}
\newcommand{\diag}{\operatorname{diag}}
\newcommand{\wh}{\widehat}
\newcommand{\norm}[1]{\left\lVert #1\right\rVert}
\newcommand{\abs}[1]{\left\lvert #1\right\rvert}
\newcommand{\eps}{\varepsilon}
\newcommand{\mcA}{\mathcal A}
\newcommand{\mcB}{\mathcal B}
\newcommand{\mcF}{\mathcal F}
\newcommand{\lesssimC}{\lesssim}

% NOTE: received as arg,min -- should be arg\,min
\DeclareMathOperator*{\argmin}{arg\,min}

\title[Matrix Spencer bounds from finite-entropy lifts]{Matrix Spencer Bounds from Finite-Entropy Lifts}
\author{Yitzchak Shmalo}
\address{Einstein Institute of Mathematics, The Hebrew University of Jerusalem, Jerusalem, Israel}
\email{yitzchak.shmalo@gmail.com}
\date{June 2026}

\subjclass[2020]{Primary 15A60, 11K38; Secondary 15B05, 42A05, 46L07, 47B35, 68Q25}
\keywords{Matrix discrepancy, Matrix Spencer conjecture, finite-dimensional C*-algebras, Toeplitz matrices, Hankel matrices, group convolution, block diagonal matrices}

\begin{document}

\begin{abstract}
The Matrix Spencer conjecture asks whether every collection of self-adjoint
$n\times n$ contractions $A_1,\ldots,A_n$ admits signs
$\eps_i\in\{-1,1\}$ for which
$\|\sum_i \eps_i A_i\|_{\op}=O(\sqrt n)$. We formulate a reusable
finite-lift principle for structured instances. Let
\[
\mcB=\bigoplus_{\alpha=1}^m \M_{r_\alpha}(\C),\qquad
h(\mcB)=\max_\alpha r_\alpha,\qquad
R(\mcB)=\sum_\alpha r_\alpha,
\]
and set $\Ent(\mcB)=h(\mcB)R(\mcB)$. If a tuple of matrices has a
$K$-bounded lift to $\mcB$ through a contraction on the lifted span, then
\[
\discop(A_1,\ldots,A_n)
\le C K \sqrt{n\log\left(e+\frac{\Ent(\mcB)}{n}\right)} .
\]
This is a direct consequence of the block-diagonal Matrix Spencer theorem of
Dadush, Jiang and Reis. The same argument also covers matrix-valued
multiplier models after sampling on a finite norming set.

The framework gives deterministic Matrix Spencer bounds for scalar Toeplitz
and Hankel contractions, fixed block-size and fixed-level Toeplitz and Hankel classes,
matrices in a common finite-dimensional C*-algebra, finite-group convolution
algebras with controlled irreducible dimensions, and the usual compression and
bounded-image operations. The paper also records the correct entropy accounting
for finite sums: for direct sums one has
$\Ent(\bigoplus_\ell \mcB_\ell)= (\max_\ell h(\mcB_\ell))\sum_\ell R(\mcB_\ell)$,
not generally $\sum_\ell \Ent(\mcB_\ell)$. Finally, an elementary dimension
argument shows that any exact finite block lift of the full algebra $\M_D(\C)$
has entropy at least $D^2$, explaining why this finite-lift mechanism alone
cannot settle the square case of the full Matrix Spencer conjecture.
\end{abstract}

\maketitle

\tableofcontents

\section{Introduction}

Spencer's theorem that ``six standard deviations suffice'' is a central result
of discrepancy theory \cite{Spencer,Matousek,Chazelle}. In one standard form,
$n$ vectors in the unit cube of $\R^n$ can be signed so that their signed sum has
$\ell_\infty$ norm $O(\sqrt n)$. The Matrix Spencer conjecture asks for the
analogous statement with vectors replaced by self-adjoint matrices and
$\ell_\infty$ norm replaced by operator norm.

\begin{definition}[Matrix discrepancy]
For self-adjoint matrices $A_1,\ldots,A_n$, set
\[
\disc(A_1,\ldots,A_n)
= \min_{\eps\in\{-1,1\}^n}
\norm{\sum_{i=1}^n \eps_i A_i}_{\op} .
\]
For arbitrary complex matrices we use the same formula and write
\[
\discop(A_1,\ldots,A_n)
=\min_{\eps\in\{-1,1\}^n}
\norm{\sum_{i=1}^n\eps_i A_i}_{\op}.
\]
\end{definition}

\begin{question}[Matrix Spencer]
Is there a universal constant $C$ such that every $n$ self-adjoint matrices
$A_1,\ldots,A_n\in \M_n(\C)$ satisfying $\norm{A_i}_{\op}\le 1$ obey
\[
\disc(A_1,\ldots,A_n)\le C\sqrt n?
\]
More generally, one asks for the same bound for $d\times d$ matrices with
$d\le n$, with a constant independent of $d$.
\end{question}

The general problem remains open. Noncommutative Khintchine inequalities give
logarithmic losses \cite{LustPisier,Pisier}. Important recent advances include
the intrinsic-freeness concentration theory of Bandeira, Boedihardjo and van
Handel \cite{BBvH}; the quantum-communication route of Hopkins, Raghavendra and
Shetty \cite{HRS}; the polylogarithmic-rank theorem of Bansal, Jiang and Meka
\cite{BJM}; and the mirror-descent matrix-discrepancy framework of Dadush,
Jiang and Reis \cite{DJR}. These results do not directly cover many full-rank
structured families, such as normalized Toeplitz, Hankel, or convolution
matrices in the square regime.

The point of this paper is to isolate a simple sufficient condition that does
cover these structured classes. The condition is the existence of a bounded
finite block lift. If each input matrix is a compression, quotient, or bounded
linear image of a block-diagonal element whose total block complexity is small,
one signs the lifted block-diagonal matrices and then maps back.

\subsection{The finite-lift principle}

Let
\[
\mcB=\bigoplus_{\alpha=1}^m \M_{r_\alpha}(\C),
\qquad
\norm{Y}_{\mcB}=\max_\alpha \norm{Y_\alpha}_{\op}.
\]
We write
\[
h(\mcB)=\max_\alpha r_\alpha,\qquad
R(\mcB)=\sum_{\alpha=1}^m r_\alpha,
\qquad
\Ent(\mcB)=h(\mcB)R(\mcB).
\]
The number $\Ent(\mcB)$ is the entropy dimension seen by block-diagonal Matrix
Spencer. It is $D$ for a diagonal algebra of dimension $D$, $D^2$ for the full
matrix algebra $\M_D$, and $bD$ for a block-diagonal algebra on $\C^D$ with
blocks of size at most $b$.

\begin{theorem}[Finite block-lift Matrix Spencer]\label{thm:main}
There is a universal constant $C$ with the following property. Let
$A_1,\ldots,A_n\in \M_D(\C)$ be arbitrary complex matrices with
$\norm{A_i}_{\op}\le 1$. Suppose there are elements $Y_i\in\mcB$ and a fixed
linear map
\[
\Pi: \Span\{Y_1,\ldots,Y_n\}\longrightarrow \M_D(\C)
\]
such that
\[
\Pi(Y_i)=A_i,
\qquad
\norm{\Pi(Y)}_{\op}\le \Lambda\norm{Y}_{\mcB}
\quad (Y\in\Span\{Y_1,\ldots,Y_n\}),
\]
and
\[
\norm{Y_i}_{\mcB}\le K \qquad (1\le i\le n).
\]
Then
\[
\discop(A_1,\ldots,A_n)
\le C\Lambda K
\sqrt{n\log\left(e+\frac{\Ent(\mcB)}{n}\right)} .
\]
If the $A_i$ are self-adjoint and the $Y_i$ may be chosen self-adjoint, the
same estimate holds for $\disc(A_1,\ldots,A_n)$.
\end{theorem}

Thus every class with uniformly bounded lift loss and $\Ent(\mcB)=O(n)$
satisfies the conjectured $O(\sqrt n)$ Matrix Spencer bound. Scalar Toeplitz
matrices are the commutative case: after a bounded finite Toeplitz extension and
Fourier sampling, the lift algebra is diagonal of dimension $O(D)$. Finite group
convolution and block Toeplitz matrices are noncommutative variants in which the
finite lift has bounded block size.

The theorem has a built-in limitation. If the full algebra $\M_D(\C)$ is lifted
exactly through a finite block algebra, surjectivity forces
$\sum_\alpha r_\alpha^2\ge D^2$, hence $\Ent(\mcB)\ge D^2$. Therefore this
mechanism cannot, by itself, turn arbitrary $D\times D$ matrices in the square
regime $D=n$ into an $O(\sqrt n)$ theorem. It yields instead the usual
$\sqrt{n\log(e+D^2/n)}$ scale. The point is not that finite lifts are the only
possible method, but that this particular method succeeds exactly when the
structured class has small finite-lift entropy.

\subsection{Main deterministic consequences}

The abstract theorem gives the following consequences. Constants are universal
unless a dependence is displayed in the subscript.

\begin{enumerate}[leftmargin=2em]
\item \textbf{Scalar Toeplitz and Hankel.} Every collection of arbitrary complex
$d\times d$ Toeplitz or Hankel contractions satisfies
\[
\discop(A_1,\ldots,A_n)
\le C\sqrt{n\log\left(e+\frac{d}{n}\right)} .
\]
In particular, self-adjoint Toeplitz and Hankel matrices satisfy the Matrix
Spencer bound when $d\le n$.

\item \textbf{Fixed-block and fixed-level Toeplitz/Hankel.} For fixed matrix
block size $b$ and fixed multilevel order $r$, all $r$-level $b\times b$ block
Toeplitz and block Hankel contractions of total dimension $D$ satisfy
\[
\discop(A_1,\ldots,A_n)
\le C_{r,b}\sqrt{n\log\left(e+\frac{D}{n}\right)} .
\]

\item \textbf{Common finite-dimensional C*-algebras.} Suppose the $A_i$ lie in a
common finite-dimensional C*-algebra represented as
\[
\mcA \simeq \bigoplus_{\alpha=1}^m
\left(\M_{r_\alpha}(\C)\otimes I_{m_\alpha}\right)
\subseteq \M_D(\C).
\]
Then
\[
\discop(A_1,\ldots,A_n)
\le C\sqrt{n\log\left(e+
\frac{(\max_\alpha r_\alpha)\sum_\alpha r_\alpha}{n}
\right)} .
\]
This includes simultaneous diagonalization, common bounded-block
diagonalization, and commutative association-scheme algebras.

\item \textbf{Finite group convolution.} For left convolution operators on
$\ell^2(G)$, where $G$ is finite and $q(G)=\max_{\pi\in\widehat G}\dim\pi$,
\[
\discop(L_{f_1},\ldots,L_{f_n})
\le C\sqrt{n\log\left(e+
\frac{q(G)\sum_{\pi\in\widehat G}\dim\pi}{n}
\right)} .
\]
For finite abelian groups this becomes
$C\sqrt{n\log(e+|G|/n)}$.

\item \textbf{Closure operations.} The framework is stable under compression,
bounded linear images, completely bounded fixed-size tensor amplification, and
finite sums when the combined direct-sum entropy remains controlled. The last
condition is important: the entropy of a direct sum is not the sum of the
individual entropies in general.

\item \textbf{Structured random inputs.} All of the preceding statements are
pathwise deterministic. They therefore apply to random Toeplitz, Hankel,
fixed-block multilevel Toeplitz/Hankel, circulant, and group-convolution samples
after normalization by their realized operator norms.
\end{enumerate}

\section{Discrepancy inputs}

We use two discrepancy inputs. The first is the rectangular vector Spencer
theorem. All logarithms are natural.

\begin{theorem}[Vector Spencer]\label{thm:vspencer}
There is a universal constant $C_S$ such that, for every
$x_1,\ldots,x_n\in\C^m$ with $\norm{x_i}_\infty\le 1$, there are signs
$\eps\in\{-1,1\}^n$ satisfying
\[
\norm{\sum_{i=1}^n \eps_i x_i}_\infty
\le C_S \sqrt{n\log\left(e+\frac{m}{n}\right)} .
\]
\end{theorem}

The complex case follows from the real theorem by applying it to real and
imaginary parts and absorbing the harmless factor of two into the constant.

The second input is the block-diagonal theorem of Dadush, Jiang and Reis. We
state it in the standard form with total matrix dimension $R$ and maximum block
size $h$.

\begin{theorem}[Block-diagonal Matrix Spencer, Dadush--Jiang--Reis]\label{thm:djr}
There is a universal constant $C_{BD}$ such that the following holds. Let
$\C^R=\bigoplus_\alpha \C^{r_\alpha}$ be a fixed block decomposition with
$h=\max_\alpha r_\alpha$. Let $Y_1,\ldots,Y_n$ be self-adjoint matrices that
are block diagonal with respect to this decomposition and satisfy
$\norm{Y_i}_{\op}\le 1$. Then there are signs $\eps\in\{-1,1\}^n$ such that
\[
\norm{\sum_{i=1}^n \eps_i Y_i}_{\op}
\le C_{BD}\sqrt{n\log\left(e+\frac{hR}{n}\right)} .
\]
Equivalently, for $\mcB=\bigoplus_\alpha \M_{r_\alpha}(\C)$,
\[
\norm{\sum_{i=1}^n \eps_i Y_i}_{\mcB}
\le C_{BD}\sqrt{n\log\left(e+\frac{\Ent(\mcB)}{n}\right)} .
\]
\end{theorem}

\begin{remark}
Dadush, Jiang and Reis prove this as an application of their partial-coloring
framework; in their notation the logarithmic factor is of the form
$\log(hm/n)$ for $m\times m$ block-diagonal matrices with block size at most
$h$ \cite{DJR}. The displayed version inserts $e+$ to cover all regimes and
uses $R$ for the total matrix dimension.
\end{remark}

For non-self-adjoint block-diagonal matrices we use Hermitian dilation. For a
matrix $Y$, set
\[
\mathscr H(Y)=\begin{pmatrix}0&Y\\Y^*&0\end{pmatrix} .
\]
Then $\mathscr H(Y)$ is self-adjoint and $\norm{\mathscr H(Y)}_{\op}=\norm{Y}_{\op}$.
Applying Theorem \ref{thm:djr} to the dilations doubles all block sizes
and changes only the universal constant.

\begin{corollary}[Non-self-adjoint block diagonal signing]\label{cor:nsa}
Let $Y_1,\ldots,Y_n\in\mcB$ be arbitrary complex block-diagonal matrices with
$\norm{Y_i}_{\mcB}\le 1$. Then
\[
\min_{\eps\in\{-1,1\}^n}
\norm{\sum_{i=1}^n \eps_iY_i}_{\mcB}
\le C\sqrt{n\log\left(e+\frac{\Ent(\mcB)}{n}\right)} .
\]
\end{corollary}

\section{Finite block and sampled multiplier lifts}

The lift formulation is tuple-level: for a given input collection it is enough
to choose bounded lifts of the individual matrices and a contraction on their
span.

\begin{definition}[Finite block lift]
Let $A_1,\ldots,A_n\in\M_D(\C)$ and
$\mcB=\bigoplus_{\alpha=1}^m \M_{r_\alpha}(\C)$. We say that the tuple admits a
$(K,\Lambda)$-lift to $\mcB$ if there exist elements $Y_i\in\mcB$ and a fixed
linear map
\[
\Pi:\Span\{Y_1,\ldots,Y_n\}\to\M_D(\C)
\]
such that
\[
\Pi(Y_i)=A_i,
\qquad
\norm{\Pi(Y)}_{\op}\le \Lambda\norm{Y}_{\mcB}
\quad (Y\in\Span\{Y_1,\ldots,Y_n\}),
\]
and $\norm{Y_i}_{\mcB}\le K$ for all $i$.
\end{definition}

The map $\Pi$ is usually a compression, a quotient map, or a fixed bounded
linear image. Contractive lifts correspond to $\Lambda=1$.

\begin{proof}[Proof of Theorem \ref{thm:main}]
First assume that the $Y_i$ are self-adjoint. Apply Theorem
\ref{thm:djr} to $Y_i/K$. There are signs for which
\[
\norm{\sum_i \eps_i Y_i}_{\mcB}
\le C K\sqrt{n\log\left(e+\frac{\Ent(\mcB)}{n}\right)} .
\]
Since $\Pi$ is linear and has norm at most $\Lambda$ on the relevant span,
\[
\norm{\sum_i \eps_i A_i}_{\op}
= \norm{\Pi\left(\sum_i \eps_iY_i\right)}_{\op}
\le \Lambda\norm{\sum_i \eps_iY_i}_{\mcB} .
\]
This proves the self-adjoint case. In the arbitrary complex case, replace each
block $Y_{i,\alpha}$ by its Hermitian dilation. The block sizes and total
reduced dimension are both doubled, so the entropy changes by a factor of four,
which is absorbed into the universal constant. The identity
$\norm{\sum_i\eps_i\mathscr H(Y_i)}=\norm{\sum_i\eps_iY_i}$ gives the stated
non-self-adjoint bound.
\end{proof}

\subsection{Sampled multiplier lifts}

Toeplitz arguments naturally produce continuous multiplier lifts rather than
finite block lifts at the first step. The following sampling principle converts
them into finite block lifts.

\begin{definition}[Norming set]
Let $\Omega$ be a compact space and let
$\mcF\subset C(\Omega;\M_b(\C))$ be a finite-dimensional linear space. A finite
set $Q\subset\Omega$ is $L_Q$-norming for $\mcF$ if
\[
\norm{F}_{L^\infty(\Omega;\M_b)}
\le L_Q\max_{\omega\in Q}\norm{F(\omega)}_{\op}
\qquad(F\in\mcF).
\]
\end{definition}

\begin{definition}[Sampled multiplier lift]
A tuple $A_1,\ldots,A_n\in\M_D(\C)$ admits a
$(K,\Lambda,L_Q,Q,b)$ sampled multiplier lift if there are functions
$F_i\in\mcF\subset C(\Omega;\M_b(\C))$ and a fixed linear map
\[
\Pi:\Span\{F_1,\ldots,F_n\}\to\M_D(\C)
\]
such that
\[
\Pi(F_i)=A_i,
\qquad
\norm{\Pi(F)}_{\op}\le \Lambda\norm{F}_{L^\infty(\Omega;\M_b)},
\]
$\norm{F_i}_{L^\infty}\le K$, and $Q$ is $L_Q$-norming for
$\Span\{F_1,\ldots,F_n\}$.
\end{definition}

\begin{theorem}[Sampled multiplier Matrix Spencer]\label{thm:sampled}
If $A_1,\ldots,A_n$ admit a $(K,\Lambda,L_Q,Q,b)$ sampled multiplier lift, then
\[
\discop(A_1,\ldots,A_n)
\le C K\Lambda L_Q
\sqrt{n\log\left(e+\frac{b^2|Q|}{n}\right)} .
\]
For $b=1$ this is the vector-Spencer scale
$CK\Lambda L_Q\sqrt{n\log(e+|Q|/n)}$.
\end{theorem}

\begin{proof}
Evaluate each $F_i$ on $Q$ and form
\[
Y_i=\bigoplus_{\omega\in Q}F_i(\omega)
\in \bigoplus_{\omega\in Q}\M_b(\C).
\]
This finite block algebra has $h=b$, $R=b|Q|$, and entropy $b^2|Q|$. By
Corollary \ref{cor:nsa}, there are signs such that
\[
\max_{\omega\in Q}
\norm{\sum_i\eps_iF_i(\omega)}_{\op}
\le C K\sqrt{n\log\left(e+\frac{b^2|Q|}{n}\right)} .
\]
Since $Q$ is $L_Q$-norming,
\[
\norm{\sum_i\eps_iF_i}_{L^\infty}
\le C K L_Q\sqrt{n\log\left(e+\frac{b^2|Q|}{n}\right)} .
\]
Applying $\Pi$ gives the result.
\end{proof}

\subsection{A grid lemma for trigonometric polynomial lifts}

We shall use the following elementary norming lemma. It is stated for
matrix-valued trigonometric polynomials because this is the form needed for
block Toeplitz matrices.

\begin{lemma}[Grid norming]\label{lem:grid}
Fix $r\ge 1$. Let $P:\T^r\to\M_b(\C)$ be a matrix-valued trigonometric
polynomial with coordinate degrees at most $N_1,\ldots,N_r$. For each $j$, let
$G_j$ be a uniform grid in $\T$ of size at least $8rN_j$ when $N_j\ge 1$, and of
size $1$ when $N_j=0$. Put $G=G_1\times\cdots\times G_r$. Then
\[
\norm{P}_{L^\infty(\T^r;\M_b)}
\le 2\max_{\theta\in G}\norm{P(\theta)}_{\op} .
\]
In particular, $|G|=O_r(\prod_{j=1}^r(1+N_j))$.
\end{lemma}

\begin{proof}
The Banach-valued Bernstein inequality gives
$\norm{\partial_j P}_\infty\le N_j\norm{P}_\infty$ for each coordinate. One way
to see this is to apply the scalar Bernstein inequality to $\varphi(P(\cdot))$
for norm-one linear functionals $\varphi$ on $\M_b$. Choose $\theta^*$ with
$\norm{P(\theta^*)}\ge\norm{P}_\infty-\delta$, and choose $\theta\in G$ with
$|\theta_j-\theta^*_j|\le \pi/|G_j|$ for every $j$. Then
\[
\norm{P(\theta^*)-P(\theta)}
\le \sum_{j=1}^r \frac{\pi}{|G_j|}N_j\norm{P}_\infty
\le \frac{\pi}{8}\norm{P}_\infty .
\]
Letting $\delta\downarrow 0$ gives
$\norm{P}_\infty\le (1-\pi/8)^{-1}\max_G\norm{P}\le 2\max_G\norm{P}$.
\end{proof}

\section{Finite-dimensional C*-algebras and group convolution}

The cleanest applications are matrices that already lie in a common
finite-dimensional C*-algebra.

\begin{theorem}[Common finite-dimensional C*-algebra]\label{thm:cstar}
Let $\mcA\subset\M_D(\C)$ be a finite-dimensional C*-algebra. Write its
representation, after unitary conjugacy, as
\[
\mcA\simeq \bigoplus_{\alpha=1}^m
\left(\M_{r_\alpha}(\C)\otimes I_{m_\alpha}\right).
\]
Let $h=\max_\alpha r_\alpha$ and $R=\sum_\alpha r_\alpha$. If
$A_1,\ldots,A_n\in\mcA$ satisfy $\norm{A_i}_{\op}\le 1$, then
\[
\discop(A_1,\ldots,A_n)
\le C\sqrt{n\log\left(e+\frac{hR}{n}\right)} .
\]
If the matrices are self-adjoint, this is a self-adjoint discrepancy bound.
\end{theorem}

\begin{proof}
An element of $\mcA$ is unitarily equivalent to
$\bigoplus_\alpha(B_\alpha\otimes I_{m_\alpha})$. Delete the multiplicity
identities to obtain the reduced algebra
$\mcB=\bigoplus_\alpha\M_{r_\alpha}(\C)$. The deletion map and the inverse map
$\bigoplus B_\alpha\mapsto\bigoplus(B_\alpha\otimes I_{m_\alpha})$ are isometric
for the operator norm. Theorem \ref{thm:main} applies with
$K=\Lambda=1$.
\end{proof}

\begin{corollary}[Simultaneous diagonalization]
If $A_1,\ldots,A_n$ are normal matrices diagonalized by a common unitary and
$\norm{A_i}_{\op}\le 1$, then
\[
\discop(A_1,\ldots,A_n)
\le C\sqrt{n\log\left(e+\frac{D}{n}\right)} .
\]
Thus the conjectured $O(\sqrt n)$ bound holds when $D\le n$.
\end{corollary}

\begin{corollary}[Common bounded-block diagonalization]
Suppose all $A_i$ are simultaneously unitarily block diagonal with common block
sizes at most $b$. Let $D$ be the ambient dimension. Then
\[
\discop(A_1,\ldots,A_n)
\le C\sqrt{n\log\left(e+\frac{bD}{n}\right)} .
\]
For fixed $b$ and $D\le n$, this is $O_b(\sqrt n)$.
\end{corollary}

\begin{corollary}[Commutative association schemes]
Let $\mcA\subset\M_D(\C)$ be a commutative self-adjoint matrix algebra with $s$
minimal spectral projections. If $A_1,\ldots,A_n\in\mcA$ are contractions, then
\[
\discop(A_1,\ldots,A_n)
\le C\sqrt{n\log\left(e+\frac{s}{n}\right)} .
\]
In particular, any family contained in the Bose--Mesner algebra of a symmetric
association scheme with $s$ primitive idempotents satisfies the Matrix Spencer
bound whenever $s=O(n)$, independently of the number of vertices.
\end{corollary}

\subsection{Finite group convolution}

Let $G$ be a finite group. A left convolution operator on $\ell^2(G)$ has the
form
\[
L_f\xi(x)=\sum_{y\in G} f(xy^{-1})\xi(y).
\]
By the finite Fourier transform,
\[
L_f \simeq \bigoplus_{\pi\in\widehat G}
\left(\widehat f(\pi)\otimes I_{d_\pi}\right),
\qquad d_\pi=\dim\pi,
\]
and
$\norm{L_f}_{\op}=\max_{\pi\in\widehat G}\norm{\widehat f(\pi)}_{\op}$.
Thus the reduced convolution algebra is
$\bigoplus_{\pi\in\widehat G}\M_{d_\pi}(\C)$.

\begin{theorem}[Finite group convolution]\label{thm:group}
Let $G$ be finite, let $q(G)=\max_{\pi\in\widehat G}d_\pi$, and put
$S(G)=\sum_{\pi\in\widehat G}d_\pi$. If
$L_{f_1},\ldots,L_{f_n}$ are left convolution operators on $\ell^2(G)$ with
$\norm{L_{f_i}}_{\op}\le 1$, then
\[
\discop(L_{f_1},\ldots,L_{f_n})
\le C\sqrt{n\log\left(e+\frac{q(G)S(G)}{n}\right)} .
\]
Consequently, for finite abelian groups,
\[
\discop(L_{f_1},\ldots,L_{f_n})
\le C\sqrt{n\log\left(e+\frac{|G|}{n}\right)} .
\]
For any family of groups with $q(G)\le q$ fixed, arbitrary convolution kernels
satisfy $O_q(\sqrt n)$ whenever $|G|\le n$.
\end{theorem}

\begin{proof}
Apply Theorem \ref{thm:cstar} to the group convolution algebra. In the
abelian case all irreducible representations have dimension one and
$S(G)=|G|$. In general $S(G)\le |G|$, and therefore $q(G)S(G)\le q|G|$ on
bounded-irrep families.
\end{proof}

\begin{remark}
Theorem \ref{thm:group} signs arbitrary convolution kernels. This is
different from the finite-group regular-representation theorem of Bandeira and
B\"olcskei \cite{BandeiraBolcskei}, which signs the specific list of left-regular
representation matrices indexed by group elements with a universal
$C\sqrt{|G|}$ bound for every finite group. The two results are complementary:
Theorem \ref{thm:group} handles arbitrary kernels but retains a logarithmic
dependence on the largest irreducible block unless the irreducible dimensions
are controlled.
\end{remark}

\section{Toeplitz, Hankel, and multilevel Toeplitz classes}

We now recover and extend the Toeplitz and Hankel endpoint. The full-bandwidth
scalar result uses bounded finite Toeplitz extension. The banded result uses an
elementary finite-section estimate and is included with full details.

\subsection{Scalar Toeplitz matrices}

For a finite sequence $(a_k)_{-(d-1)\le k\le d-1}$, write
\[
T_d(a)= (a_{r-s})_{0\le r,s<d} .
\]
If $\varphi\in L^\infty(\T)$ has Fourier coefficients
$\widehat\varphi(k)=a_k$ for $|k|<d$, then $\varphi$ induces $T_d(a)$.

\begin{theorem}[Finite Toeplitz extension]\label{thm:text}
There is an absolute constant $K_T$ such that every complex $d\times d$ Toeplitz
matrix $A=T_d(a)$ has an inducer $\varphi_A\in L^\infty(\T)$ satisfying
\[
\widehat\varphi_A(k)=a_k \quad (|k|<d),
\qquad
\norm{\varphi_A}_{L^\infty}\le K_T\norm{A}_{\op} .
\]
One may take $K_T\le 3$.
\end{theorem}

This is the finite Toeplitz extension theorem of Nikolskaya and Farforovskaya
\cite{NikolskayaFarforovskaya}; Courtney and Sarason record the bound and note
that the same reasoning applies in the non-self-adjoint case \cite{CourtneySarason}.

We also use a standard de la Vall\'ee Poussin kernel $V_N$ satisfying
\[
\widehat V_N(k)=1\quad(|k|\le N),
\qquad
\widehat V_N(k)=0\quad(|k|>2N),
\qquad
\norm{V_N}_{L^1}\le C_V .
\]

\begin{theorem}[Arbitrary scalar Toeplitz signing]\label{thm:toep}
There is a universal constant $C$ such that, for every collection of arbitrary
complex Toeplitz matrices $A_1,\ldots,A_n\in\M_d(\C)$ with
$\norm{A_i}_{\op}\le 1$,
\[
\discop(A_1,\ldots,A_n)
\le C\sqrt{n\log\left(e+\frac{d}{n}\right)} .
\]
In particular, self-adjoint Toeplitz matrices satisfy Matrix Spencer in the
regime $d\le n$.
\end{theorem}

\begin{proof}
For each $i$, choose an inducer $\varphi_i$ from Theorem
\ref{thm:text}. Put
$p_i=\varphi_i*V_{d-1}$. Then $p_i$ is a trigonometric polynomial of degree at
most $2d-2$, its Fourier coefficients agree with those of $\varphi_i$ for
$|k|<d$, and
\[
\norm{p_i}_{\infty}\le C_VK_T .
\]
Hence $T_d(p_i)=A_i$.

Choose a uniform grid $Q\subset\T$ of size $O(d)$ that is $2$-norming for
trigonometric polynomials of degree at most $2d$ by Lemma \ref{lem:grid}.
Apply Theorem \ref{thm:sampled} with $b=1$, $K=C_VK_T$, and $\Lambda=1$. There
are signs for which the signed polynomial $P=\sum_i\eps_i p_i$ satisfies
\[
\norm{P}_{\infty}
\le C\sqrt{n\log\left(e+\frac{d}{n}\right)} .
\]
Toeplitz compression is contractive from $L^\infty(\T)$ to operator norm, so
\[
\norm{\sum_i\eps_i A_i}_{\op}
= \norm{T_d(P)}_{\op}
\le \norm{P}_{\infty} .
\]
\end{proof}

\begin{remark}[Circulant-lift interpretation]
If the grid size is larger than twice the degree of the $p_i$, then the vector
$(p_i(\theta))_{\theta\in Q}$ defines a circulant matrix whose upper-left
$d\times d$ block is $A_i$. Thus the preceding proof is a bounded
circulant-lift proof.
\end{remark}

\subsection{Scalar Hankel matrices}

Let $J_d$ be the reversal matrix, $(J_dx)_r=x_{d-1-r}$. If $H$ is Hankel, then
$HJ_d$ is Toeplitz and $\norm{HJ_d}=\norm H$.

\begin{corollary}[Arbitrary scalar Hankel signing]\label{cor:hankel}
For arbitrary complex Hankel matrices $H_1,\ldots,H_n\in\M_d(\C)$ with
$\norm{H_i}_{\op}\le 1$,
\[
\discop(H_1,\ldots,H_n)
\le C\sqrt{n\log\left(e+\frac{d}{n}\right)} .
\]
Thus self-adjoint Hankel matrices satisfy Matrix Spencer when $d\le n$.
\end{corollary}

\begin{proof}
Apply Theorem \ref{thm:toep} to the Toeplitz matrices $H_iJ_d$ and
observe that
\[
\norm{\sum_i\eps_iH_i}
=\norm{\left(\sum_i\eps_iH_iJ_d\right)J_d}
=\norm{\sum_i\eps_iH_iJ_d} .
\]
\end{proof}

\subsection{Banded Toeplitz matrices}

For bandwidth $w\ll d$, one needs only $O(w)$ Fourier samples. The following
finite-section lemma replaces the bounded extension theorem in this case.

\begin{lemma}[Finite-section lower bound for banded symbols]\label{lem:band}
Let $p(e^{it})=\sum_{|k|\le w}a_ke^{ikt}$ and let
$T_d(p)=(a_{r-s})_{0\le r,s<d}$. If $d>2w$, then
\[
\norm{p}_{\infty}
\le \left(1-\frac{2w}{d}\right)^{-1/2}\norm{T_d(p)}_{\op} .
\]
In particular, if $d\ge 4w$, then $\norm{p}_{\infty}\le \sqrt2\norm{T_d(p)}$.
\end{lemma}

\begin{proof}
Choose $t_0$ such that $|p(e^{-it_0})|=\norm p_\infty$, and set
$x_s=d^{-1/2}e^{is t_0}$ for $0\le s<d$. Then $\norm x_2=1$. For every
interior row $w\le r\le d-1-w$,
\[
(T_d(p)x)_r
= \frac{1}{\sqrt d}\sum_{|k|\le w} a_k e^{i(r-k)t_0}
= \frac{e^{ir t_0}}{\sqrt d}p(e^{-it_0}).
\]
Therefore
\[
\norm{T_d(p)}_{\op}^2
\ge \norm{T_d(p)x}_2^2
\ge \frac{d-2w}{d}\norm p_\infty^2 .
\]
\end{proof}

\begin{theorem}[Banded scalar Toeplitz signing]\label{thm:banded}
There is a universal constant $C$ such that the following holds. Let
$A_i=T_d(p_i)$ be arbitrary complex Toeplitz matrices whose symbols have degree
at most $w$, and assume $\norm{A_i}_{\op}\le 1$. If $d\ge 4w$ and $w\ge 1$,
then
\[
\discop(A_1,\ldots,A_n)
\le C\sqrt{n\log\left(e+\frac{w}{n}\right)} .
\]
For $w=0$ the same statement holds with $w$ replaced by $1$ inside the logarithm.
\end{theorem}

\begin{proof}
By Lemma \ref{lem:band}, $\norm{p_i}_\infty\le\sqrt2$ for all $i$.
Choose an $O(w)$-point uniform grid that is $2$-norming for degree-$w$
trigonometric polynomials by Lemma \ref{lem:grid}. Apply Theorem
\ref{thm:sampled} with $b=1$ to obtain signs with
\[
\norm{\sum_i\eps_i p_i}_\infty
\le C\sqrt{n\log\left(e+\frac{w}{n}\right)} .
\]
Contractivity of Toeplitz compression gives the same bound for
$\sum_i\eps_iA_i$.
\end{proof}

\subsection{Block and multilevel Toeplitz matrices}

An $r$-level $b\times b$ block Toeplitz matrix is indexed by a box
\[
\Lambda=\prod_{j=1}^r\{0,\ldots,d_j-1\}
\]
and has $b\times b$ blocks depending only on coordinate differences. Its total
dimension is
\[
D=b\prod_{j=1}^r d_j .
\]

We use the following finite multilevel Toeplitz extension theorem as a black
box. For $r=1$ this is Theorem \ref{thm:text}. For $r>1$ it is
the finite discrete form obtained from the multivariable Toeplitz/Nehari
extension theorem of Carlsson and Perfekt \cite{CarlssonPerfekt}.

% VERIFICATION FLAG: this is the load-bearing claim. The dimension-only constant
% K_r for r>=2 is exactly where multivariable Nehari/bounded-extension is delicate
% and may FAIL with comparable constants. Verify or restrict to r=1 + banded multilevel.
\begin{theorem}[Finite multilevel Toeplitz extension]\label{thm:mltext}
For every fixed $r\ge 1$ there is a constant $K_r$ such that every scalar
$r$-level Toeplitz matrix $A$ on the box $\Lambda$ has an inducer
$\varphi_A\in L^\infty(\T^r)$ satisfying
\[
\widehat\varphi_A(k)=a_k
\quad\text{for all difference indices } k\in\prod_{j=1}^r\{-(d_j-1),\ldots,d_j-1\},
\]
and
\[
\norm{\varphi_A}_{L^\infty(\T^r)}\le K_r\norm A_{\op} .
\]
\end{theorem}

\begin{theorem}[Fixed-level block Toeplitz and Hankel]\label{thm:blocktoep}
Fix $r\ge 1$ and $b\ge 1$. There is a constant $C_{r,b}$ such that every
collection of $r$-level $b\times b$ block Toeplitz contractions
$A_1,\ldots,A_n\in\M_D(\C)$ satisfies
\[
\discop(A_1,\ldots,A_n)
\le C_{r,b}\sqrt{n\log\left(e+\frac{D}{n}\right)} .
\]
The same bound holds for the corresponding $r$-level block Hankel class. Hence,
for fixed $r,b$ and $D\le n$, these classes satisfy the Matrix Spencer bound
$O_{r,b}(\sqrt n)$.
\end{theorem}

\begin{proof}
First consider block Toeplitz matrices. For a fixed pair of block coordinates
$p,q\in\{1,\ldots,b\}$, the scalar $r$-level Toeplitz matrix formed by the
$(p,q)$ entries of all blocks is a compression of $A_i$ between coordinate
subspaces, hence has norm at most $\norm{A_i}_{\op}$. By Theorem
\ref{thm:mltext}, each scalar entry has an $L^\infty(\T^r)$ inducer
with norm at most $K_r\norm{A_i}$. Combining the $b^2$ scalar inducers gives a
matrix-valued inducer $\Phi_i\in L^\infty(\T^r;\M_b)$ with
\[
\norm{\Phi_i}_{L^\infty(\T^r;\M_b)}\le bK_r\norm{A_i}_{\op} .
\]

Convolve $\Phi_i$ with the product de la Vall\'ee Poussin kernel
$\prod_{j=1}^r V_{d_j-1}$ to obtain a matrix trigonometric polynomial $P_i$.
Its coordinate degrees are $O(d_j)$, it has the same finite Toeplitz
coefficients as $A_i$, and
\[
\norm{P_i}_{L^\infty(\T^r;\M_b)}\le C_{r,b}\norm{A_i}_{\op} .
\]
By Lemma \ref{lem:grid}, there is a $2$-norming grid
$Q\subset\T^r$ of size
\[
|Q|=O_r\left(\prod_{j=1}^r d_j\right)=O_r(D/b).
\]
Theorem \ref{thm:sampled}, with block size $b$, gives
\[
\discop(A_1,\ldots,A_n)
\le C_{r,b}\sqrt{n\log\left(e+\frac{b^2|Q|}{n}\right)}
\le C_{r,b}\sqrt{n\log\left(e+\frac{D}{n}\right)},
\]
because $r$ and $b$ are fixed. The block Hankel statement follows by multiplying
by the appropriate tensor-product reversal matrix, exactly as in Corollary
\ref{cor:hankel}.
\end{proof}

\section{Closure operations}

The finite-lift framework is stable under the operations normally used to build
structured matrices. The entropy accounting in this section is explicit because
this is a common source of mistakes.

\begin{proposition}[Compressions and bounded images]\label{prop:comp}
Suppose $B_i$ satisfy a discrepancy bound
\[
\discop(B_1,\ldots,B_n)\le \Delta .
\]
If $A_i=V^*B_iV$ for a fixed isometry $V$, then
$\discop(A_1,\ldots,A_n)\le\Delta$. More generally, if
$A_i=\Phi(B_i)$ for a fixed linear map $\Phi$ with $\norm\Phi\le L$ on the
relevant span, then
\[
\discop(A_1,\ldots,A_n)
\le L\Delta .
\]
\end{proposition}

\begin{proof}
Use the same signs and the inequality
$\norm{\sum_i\eps_i\Phi(B_i)}\le L\norm{\sum_i\eps_iB_i}$.
\end{proof}

\begin{proposition}[Finite sums and direct-sum entropy]\label{prop:sums}
Suppose
\[
A_i=A_i^{(1)}+\cdots+A_i^{(s)} .
\]
For each $\ell$, assume the tuple $(A_i^{(\ell)})_{i=1}^n$ has a
$(K_\ell,\Lambda_\ell)$-lift to
\[
\mcB_\ell=\bigoplus_\alpha \M_{r_{\ell,\alpha}}(\C).
\]
Put
\[
h=\max_{\ell,\alpha} r_{\ell,\alpha},
\qquad
R=\sum_{\ell=1}^s\sum_\alpha r_{\ell,\alpha} .
\]
Then $(A_i)_{i=1}^n$ has a lift to $\bigoplus_{\ell=1}^s\mcB_\ell$ and
\[
\discop(A_1,\ldots,A_n)
\le C\left(\sum_{\ell=1}^s\Lambda_\ell K_\ell\right)
\sqrt{n\log\left(e+\frac{hR}{n}\right)} .
\]
\end{proposition}

\begin{proof}
For each component choose lifts $Y_i^{(\ell)}$ in $\mcB_\ell$. If $K_\ell>0$,
set $Z_i^{(\ell)}=Y_i^{(\ell)}/K_\ell$; zero components may be omitted. Then
$Z_i=(Z_i^{(1)},\ldots,Z_i^{(s)})$ has norm at most one in
$\bigoplus_\ell\mcB_\ell$. Define
\[
\Theta(Z^{(1)},\ldots,Z^{(s)})
=\sum_{\ell=1}^s K_\ell\Pi_\ell(Z^{(\ell)}).
\]
Then $\Theta(Z_i)=A_i$ and
$\norm{\Theta}\le\sum_\ell\Lambda_\ell K_\ell$ on the relevant span. Apply
Theorem \ref{thm:main} to the direct-sum algebra, whose maximum block
size is $h$ and whose reduced dimension is $R$.
\end{proof}

\begin{remark}[Direct sums are not additive in entropy]\label{rem:nonadd}
For $\mcB=\bigoplus_{\ell=1}^s\mcB_\ell$ one has
\[
\Ent(\mcB)=\left(\max_\ell h(\mcB_\ell)\right)
\sum_{\ell=1}^s R(\mcB_\ell),
\]
not generally $\sum_\ell\Ent(\mcB_\ell)$. Consequently, a finite sum of classes
that are individually in the Spencer regime remains in the Spencer regime only
when this combined entropy is $O(n)$. This condition is automatic, for example,
for a fixed number of scalar sampled multiplier classes of total sample size
$O(n)$, and for a fixed number of uniformly bounded-block classes with total
reduced dimension $O(n)$.
\end{remark}

\begin{corollary}[Toeplitz plus Hankel with bounded decomposition]
If $A_i=T_i+H_i$, where $T_i$ is Toeplitz, $H_i$ is Hankel, and
\[
\norm{T_i}_{\op}+\norm{H_i}_{\op}\le K
\qquad(1\le i\le n),
\]
then
\[
\discop(A_1,\ldots,A_n)
\le CK\sqrt{n\log\left(e+\frac{d}{n}\right)} .
\]
\end{corollary}

\begin{proof}
Use the scalar Toeplitz and Hankel multiplier lifts from Theorem
\ref{thm:toep} and Corollary \ref{cor:hankel}. Their sampled lift
algebras are diagonal with $O(d)$ total samples. The direct sum of the two
sampled algebras is still diagonal and has $O(d)$ total samples. Proposition
\ref{prop:sums} gives the result.
\end{proof}

\begin{proposition}[Completely bounded tensor amplification]\label{prop:cb}
Let $S\subset\M_D(\C)$ be an operator space and suppose there are linear maps
\[
\sigma:S\to\mcB=\bigoplus_\alpha\M_{r_\alpha}(\C),
\qquad
\Pi:\sigma(S)\to S,
\]
with $\Pi\sigma=\operatorname{id}_S$, $\norm{\sigma}_{cb}\le K$, and
$\norm{\Pi}_{cb}\le\Lambda$. Then, for each $c\ge1$, the amplified space
$S\otimes\M_c$ has a $(K,\Lambda)$ lift to
\[
\mcB\otimes\M_c\simeq\bigoplus_\alpha\M_{cr_\alpha}(\C).
\]
The entropy is multiplied by $c^2$. Hence fixed-size completely bounded tensor
amplification preserves the Spencer-regime conclusion.
\end{proposition}

\begin{proof}
Use the maps $\sigma\otimes\operatorname{id}_{\M_c}$ and
$\Pi\otimes\operatorname{id}_{\M_c}$. The cb-norm assumptions give the same
lift constants at matrix level $c$. The new block sizes are $cr_\alpha$, so both
$h$ and $R$ are multiplied by $c$.
\end{proof}

\begin{remark}
The completely bounded hypothesis in Proposition \ref{prop:cb} is necessary
for this general formulation. For concrete lifts coming from $*$-homomorphisms,
compressions, or C*-quotient maps it is automatically satisfied with the
expected constants.
\end{remark}

\section{Random structured inputs}

The preceding results are deterministic. They therefore apply pathwise to
random samples from the covered structured classes after normalization.

\begin{corollary}[Pathwise random structured matrices]
Let $A_1,\ldots,A_n$ be random matrices which, almost surely, lie in one of the
classes covered above: scalar Toeplitz, scalar Hankel, fixed-level block
Toeplitz/Hankel, a fixed finite-dimensional C*-algebra of entropy $\Ent$, or a
finite group convolution algebra. On each realization define
\[
\widetilde A_i=
\begin{cases}
A_i/\norm{A_i}_{\op},& A_i\ne0,\\
0,& A_i=0.
\end{cases}
\]
Then the corresponding deterministic bound holds for
$\widetilde A_1,\ldots,\widetilde A_n$ on that realization.
\end{corollary}

\begin{remark}
This pathwise statement is not a deterministic Matrix Spencer theorem for
generic random matrix realizations. GOE, Wigner, and Wishart matrices typically
span high-entropy subspaces of the full matrix algebra, and the finite-lift
method here does not remove the full-algebra logarithmic loss for such generic
inputs. Average-case matrix discrepancy for these ensembles is a separate line
of work; see, for example, Kunisky--Zhang \cite{KuniskyZhang}, Maillard
\cite{Maillard}, and the later combined work of Kunisky, Oertel, Wengiel and
Zhang \cite{KOWZ}.
\end{remark}

\section{Boundary of the finite-lift method}

The framework gives a precise limitation for exact finite block lifts.

\begin{proposition}[Entropy obstruction for the full matrix algebra]\label{prop:obstruct}
Suppose the full algebra $\M_D(\C)$ admits an exact finite block lift through
$\mcB=\bigoplus_\alpha\M_{r_\alpha}(\C)$ in the sense that there is a linear map
from $\mcB$ onto $\M_D(\C)$. Then
\[
\Ent(\mcB)
=\left(\max_\alpha r_\alpha\right)\sum_\alpha r_\alpha
\ge D^2 .
\]
Consequently, Theorem \ref{thm:main} cannot yield better than
$C\sqrt{n\log(e+D^2/n)}$ for arbitrary $D\times D$ matrices through exact finite
block lifts.
\end{proposition}

\begin{proof}
Surjectivity onto the $D^2$-dimensional vector space $\M_D(\C)$ forces
\[
\dim\mcB=\sum_\alpha r_\alpha^2\ge D^2 .
\]
Since
\[
\sum_\alpha r_\alpha^2
\le \left(\max_\alpha r_\alpha\right)\sum_\alpha r_\alpha
=\Ent(\mcB),
\]
the claim follows.
\end{proof}

This explains why Toeplitz and Hankel classes are tractable by this method:
the relevant finite Fourier data has only $O(D)$ scalar degrees of freedom and
admits a bounded $L^\infty$ extension. It also explains why common block
diagonalization with bounded blocks is tractable: the reduced noncommutative
entropy is linear in $D$. Conversely, the full matrix algebra has entropy
$D^2$.

\begin{question}[Bounded finite-entropy lift problem]
Classify structured operator systems $S_D\subset\M_D(\C)$ for which every
contraction in $S_D$ has a uniformly bounded lift either to a finite block
algebra $\bigoplus_\alpha\M_{r_\alpha}(\C)$ with
\[
\left(\max_\alpha r_\alpha\right)\sum_\alpha r_\alpha=O(D),
\]
or to a matrix-valued multiplier model with an $O(D)$-size norming set and
bounded block size. Every such class satisfies Matrix Spencer in the regime
$D\le n$.
\end{question}

The Toeplitz, Hankel, fixed-level multilevel Toeplitz, finite abelian
convolution, bounded-irrep convolution, simultaneous diagonalization,
bounded-block diagonalization, and commutative association-scheme examples are
instances of this program. Solving the full Matrix Spencer conjecture will
require either a way to beat the entropy obstruction for generic matrix algebras
or a method not reducible to exact finite block lifts.

\section*{Acknowledgements and disclosure}
The author is responsible for all mathematical claims and bibliographic
information. Generative AI tools were used to accelerate drafting and revision;
all arguments and references should be independently checked before submission.

\begin{thebibliography}{99}

\bibitem{BBvH}
A.~S. Bandeira, M.~T. Boedihardjo, and R. van Handel,
\emph{Matrix concentration inequalities and free probability},
Invent. Math. \textbf{234} (2023), 419--487.

\bibitem{BandeiraBolcskei}
A.~S. Bandeira and H. B\"olcskei,
\emph{Matrix discrepancy for representations of finite groups},
arXiv:2606.12181, 2026.

\bibitem{BJM}
N. Bansal, H. Jiang, and R. Meka,
\emph{Resolving Matrix Spencer conjecture up to poly-logarithmic rank},
SIAM J. Comput., Special Section STOC 2023,
DOI: 10.1137/23M1592201.

\bibitem{BottcherGrudsky}
A. B\"ottcher and S.~M. Grudsky,
\emph{Spectral Properties of Banded Toeplitz Matrices},
SIAM, Philadelphia, 2005.

\bibitem{CarlssonPerfekt}
M. Carlsson and K.-M. Perfekt,
\emph{Nehari's theorem for convex domain Hankel and Toeplitz operators},
Int. Math. Res. Not. IMRN \textbf{2021}, no.~5, 3331--3371.

\bibitem{Chazelle}
B. Chazelle,
\emph{The Discrepancy Method: Randomness and Complexity},
Cambridge University Press, 2000.

\bibitem{CourtneySarason}
D. Courtney and D. Sarason,
\emph{A mini-max problem for self-adjoint Toeplitz matrices},
Math. Scand. \textbf{110} (2012), 82--98.

\bibitem{DJR}
D. Dadush, H. Jiang, and V. Reis,
\emph{A new framework for matrix discrepancy: partial coloring bounds via mirror descent},
Proceedings of the 54th Annual ACM SIGACT Symposium on Theory of Computing
(STOC), 2022, 649--658; arXiv:2111.03171.

\bibitem{Gluskin}
E.~D. Gluskin,
\emph{Extremal properties of orthogonal parallelepipeds and their applications to the geometry of Banach spaces},
Math. USSR-Sb. \textbf{64} (1989), 85--96.

\bibitem{HRS}
S.~B. Hopkins, P. Raghavendra, and A. Shetty,
\emph{Matrix discrepancy from quantum communication},
Proceedings of the 54th Annual ACM SIGACT Symposium on Theory of Computing
(STOC), 2022, 637--648.

\bibitem{KuniskyZhang}
D. Kunisky and P. Zhang,
\emph{Average-case matrix discrepancy: asymptotics and online algorithms},
arXiv:2307.10055, 2024.

\bibitem{LustPisier}
F. Lust-Piquard and G. Pisier,
\emph{Non-commutative Khintchine and Paley inequalities},
Ark. Mat. \textbf{29} (1991), 241--260.

\bibitem{Maillard}
A. Maillard,
\emph{Average-case matrix discrepancy: satisfiability bounds},
arXiv:2410.17887, 2024.

\bibitem{Matousek}
J. Matou\v{s}ek,
\emph{Geometric Discrepancy: An Illustrated Guide},
Springer, 1999.

\bibitem{NikolskayaFarforovskaya}
L.~N. Nikolskaya and Yu.~B. Farforovskaya,
\emph{Toeplitz and Hankel matrices as Hadamard--Schur multipliers},
Algebra i Analiz \textbf{15} (2003), no.~6, 141--160.

\bibitem{Pisier}
G. Pisier,
\emph{Introduction to Operator Space Theory},
Cambridge University Press, 2003.

\bibitem{Spencer}
J. Spencer,
\emph{Six standard deviations suffice},
Trans. Amer. Math. Soc. \textbf{289} (1985), 679--706.

\bibitem{KOWZ}
D. Kunisky, T. Oertel, N. Wengiel, and P. Zhang,
\emph{Asymptotic bounds and online algorithms for average-case matrix discrepancy},
arXiv:2410.23915v2, 2025.

\end{thebibliography}

\end{document}
